\documentclass[a4paper,fleqn]{cas-sc}

\usepackage{float}
\usepackage[utf8]{inputenc}
\usepackage[numbers]{natbib}
\usepackage{tabularx}
\usepackage{array}
\usepackage[english]{babel}
\usepackage{amsmath}
\usepackage{amsfonts}
\usepackage{amssymb}
\usepackage{booktabs}
% Float placement: the manuscript carries 7 tables and 4 figures. Allow pages
% to hold more than two floats and keep text pages from breaking early.
\setcounter{topnumber}{3}
\setcounter{bottomnumber}{2}
\setcounter{totalnumber}{5}
\renewcommand{\topfraction}{0.95}
\renewcommand{\bottomfraction}{0.95}
\renewcommand{\textfraction}{0.05}
\renewcommand{\floatpagefraction}{0.7}
% Tighter float and caption spacing, set at begin-document so the class
% defaults cannot override them.
\AtBeginDocument{%
  \setlength{\textfloatsep}{10pt plus 2pt minus 2pt}%
  \setlength{\floatsep}{8pt plus 2pt minus 2pt}%
  \setlength{\intextsep}{8pt plus 2pt minus 2pt}%
  \setlength{\abovecaptionskip}{6pt}%
  \setlength{\belowcaptionskip}{2pt}%
}
\usepackage{multirow}
\usepackage{hyperref}
\usepackage{color}
\usepackage{xcolor}
\usepackage{orcidlink}
\usepackage{lineno}
\usepackage{silence}
\WarningFilter{hyperref}{Ignoring empty anchor}
\setlength{\emergencystretch}{2em}
\setlength{\textfloatsep}{4pt plus 2pt minus 2pt}
\setlength{\floatsep}{4pt plus 2pt minus 2pt}
\setlength{\intextsep}{4pt plus 2pt minus 2pt}
\setlength{\bibsep}{0.2pt plus 0.1pt minus 0.1pt}
% Table styling: plain grid tables with roomier cells and rows.
\renewcommand{\arraystretch}{1.08}
\setlength{\tabcolsep}{5pt}

% Manuscript-specific markers. Keep them visible: an unmarked placeholder is
% a fabrication risk. \pending{...} is the only permitted placeholder for a
% confirmatory quantity.
\newcommand{\pending}[1]{\textcolor{red}{[PENDING: #1]}}
\newcommand{\tbd}{\textcolor{red}{[TBD]}}

% Figure helper: include the file when it exists, otherwise draw a labelled
% placeholder box. The four figures are generated by paper/scripts/make_figures.py;
% the manuscript must stay readable (and compilable) before they land.
\newcommand{\figfile}[2][\linewidth]{%
  \IfFileExists{#2}%
    {\includegraphics[width=#1]{#2}}%
    {\fbox{\begin{minipage}[c][5.4cm][c]{0.90\linewidth}%
      \centering\ttfamily\small #2\\[4pt]%
      \normalfont\itshape figure pending\end{minipage}}}%
}

% Preprint footer: drop the "submitted to Elsevier" wording.
\ExplSyntaxOn
\cs_set:Npn \__first_footerline:
  {
    \group_begin:
    \small
    \sffamily
    \ifnum\theblind>0\relax
    \else
    \__short_authors: :~
    \fi
    \rmfamily \itshape Preprint
    \group_end:
  }
\ExplSyntaxOff
% The CAS class draws email/URL icons from a thumbnails/ folder. Uploads that
% do not carry that folder fail with "File `thumbnails/cas-email.jpeg' not
% found". The class has a no-logo mode, so the email footnote prints as plain
% text and the document needs no image assets.
\ExplSyntaxOn
\bool_gset_true:N \g_stm_nologo_bool
\ExplSyntaxOff
% --- ORCID badges in the byline -------------------------------------------
% The CAS class prints ORCID iDs only in the footnote. Below is the class's
% own \author definition (cas-common.sty v2.3) with one addition: when an
% author declares [orcid=...], the green iD badge follows the name before the
% affiliation mark. Every other line is unchanged from the class.
\makeatletter
\ExplSyntaxOn
\RenewDocumentCommand \author { O{} m O{} }
   {
     \ResetMarks
     \tl_if_blank:nTF { #3 } { }
       { \keys_set:nn { stm / author } { #3 } }
%
     \int_gincr:N \g_stm_au_int
     \tex_gdef:D \theau@ { \int_use:N \g_stm_au_int }
%
     \seq_gput_right:Nn \g_stm_prelimsau_seq { #2 }
     \bool_if:NTF \l_stm_au_collab_bool
       { \seq_gput_right:cn { g_stm_clau\int_use:N \g_stm_augr_int _seq } }
       { \seq_gput_right:cn { g_stm_au\int_use:N \g_stm_augr_int _seq } }
       {
         \int_gincr:N \g_stm_aau_int
         \tex_gdef:D \theauthor {\int_use:N \g_stm_aau_int }
         \keys_set:nn { stm /author } { #3 }
         \tl_if_head_eq_catcode:nNTF { #1 } a
            { \processAffRef { #1 } }
            { \processAffNum { #1 } }
         \tl_if_empty:NF \l_stm_au_prefix
            { \l_stm_au_prefix_tl \c_space_token }
         \str_if_eq:VnTF \l_stm_au_style_tl  { chinese }
           {
             \invparsename { #2 }
             \textcolor{\l_stm_augroup_color_tl}{\surname}
             \no_break_space:
             \textcolor{\l_stm_augroup_color_tl !50}{\firstname}
           }
           {
             \parsename { #2 }
             \textcolor{\l_stm_augroup_color_tl !50}{\firstname}
             \no_break_space:
             \textcolor{\l_stm_augroup_color_tl}{\surname}
           }
       \tl_if_empty:NF \l_stm_au_suffix_tl
       { \c_space_token \l_stm_au_suffix_tl }
        \tl_if_empty:NF \l_stm_au_orcid_tl
          { \,\small \exp_args:NV \orcidlink \l_stm_au_orcid_tl }
        \unskip
        \textsuperscript
        {
          \tl_if_blank:nTF { #1 }
          { \tex_def:D \sep{} }
          { {\itshape\listAff} \tex_def:D \sep{\unskip,} }
          \process@marks
          \bool_if:NT \l_stm_au_deceased_bool
          { \sep \maltese
            \tex_def:D \sep { \unksip, }
          }
        }
       \tl_if_empty:NF \l_stm_au_degree_tl
       { ,\c_space_token \l_stm_au_degree_tl }
       \tl_if_empty:NF \l_stm_au_role_tl
       { \c_space_token (\l_stm_au_role_tl) }
%
      \ResetMarks
     }
%
     \bool_if:NT \l_stm_au_deceased_bool
      {
       \seq_gput_right:Nn \g_stm_maltese_seq
        {
         \tex_def:D \thefootnote { \maltese }
         \footnotetext{Deceased~author.}
        }
      }
% various social media
     \tl_if_empty:NTF \l_stm_au_facebook_tl { }
       {
         \wrAux { \token_to_str:N \facebookauthor
           { \l_stm_au_facebook_tl } { \exp_not:n {#2} } }
       }
     \tl_if_empty:NF \l_stm_au_twitter_tl
       {
         \wrAux { \token_to_str:N \twitterauthor
           { \l_stm_au_twitter_tl } { \exp_not:n {#2} } }
       }
     \tl_if_empty:NF \l_stm_au_gplus_tl
       {
         \wrAux { \token_to_str:N \gplusauthor
           { \l_stm_au_gplus_tl } { \exp_not:n {#2} } }
       }
     \tl_if_empty:NF \l_stm_au_linkedin_tl
       {
         \wrAux { \token_to_str:N \linkedinauthor
           { \l_stm_au_linkedin_tl } { \exp_not:n {#2} } }
       }
     \tl_if_empty:NF \l_stm_au_orcid_tl
       {
         \wrAux { \token_to_str:N \orcidauthor
           { \l_stm_au_orcid_tl } { \exp_not:n {#2} } }
       }
%
       \@eadauthor={#2}
       \pdfstringdef\__info_au: { #2 }
       \int_compare:nNnTF { \theau@ } < { 4 }
       { \xappto \infoauthors { \__info_au: , ~ } }
       {
         \int_compare:nNnTF { \theau@ } = { 4 }
         { \xappto \infoauthors { et~al. } }
         { }
       }
   }
\ExplSyntaxOff
\makeatother
% The byline badges already carry the iD links, so the class's separate
% "ORCID(s):" footnote would only repeat the same numbers.
\RenewDocumentCommand \printorcid { } { }
% Same author list printer as the class, except that the last two names are
% separated by a comma instead of "and" so the four-author line fits on one
% line together with the ORCID badges.
\ExplSyntaxOn
\RenewDocumentCommand \stmauthors { }
   {
     \group_begin:
     \stmAuthorSetup { type = authors }
     \l_stm_au_setup_tl
     \par \vskip\l_stm_augroup_before_dim
     \l_stm_augroup_align_tl
     \l_stm_augroup_size_tl
     \l_stm_augroup_shape_tl
     \l_stm_augroup_weight_tl
     \color{ \l_stm_augroup_color_tl }
     \bool_if:NTF \g_stm_augr_bool
       { \seq_use:Nn \g_stm_augr_seq { \par } }
       {
         \seq_use:cnnn { g_stm_au\int_use:N\g_stm_augr_int _seq }
         { ,~ } { ,~ } { ,~ }
       }
     \par\vskip\l_stm_augroup_after_dim
     \group_end:
   }
\ExplSyntaxOff

% Figures live in paper/figures/.
\graphicspath{{./figures/}{figures/}{paper/figures/}}

\begin{document}
\let\WriteBookmarks\relax

% Short title
\shorttitle{Tool-Result Errors in Language-Model Agents}

% Short author
\shortauthors{Rahaman et~al.}

% Main title of the paper
\title [mode = title]{Tool-Result Errors in Language-Model Agents: A Fixed-Prefix Causal Study of Detection, Reliance, and Recovery}

% ---------------------------------------------------------------------------
% Author block. The group, affiliations, ORCID iDs, and CRediT roles are
% carried over from the companion manuscript. CONFIRM WITH ALL AUTHORS before
% submission: author order, spelling, e-mail addresses, ORCID iDs, and the
% CRediT allocation below must be checked and signed off.
% ---------------------------------------------------------------------------

% First author and corresponding author
\author[1]{Md. Atikur Rahaman}[orcid=0009-0008-2945-4117]
\cormark[1]
\ead{mrahaman2310298@bscse.uiu.ac.bd}
\credit{Conceptualization, Data curation, Formal analysis, Investigation, Methodology, Project administration, Resources, Software, Validation, Visualization, Writing - original draft, Writing - review \& editing}



% Supervisor
\author[1]{Ohidujjaman Tuhin}[orcid=0000-0002-8776-7145]
\credit{Supervision, Validation, Writing - review \& editing}

% Address/affiliation
\affiliation[1]{organization={Computer Science and Engineering, United International University},
  city={Dhaka},
  citysep={},
  postcode={1212},
  country={Bangladesh}}

\cortext[cor1]{Corresponding author: Md. Atikur Rahaman}

% Abstract. Design numbers are sourced from README.md, ROADMAP.md,
% docs/PROTOCOL.md, and docs/DATA_CARD.md. Confirmatory values do not exist.
\begin{abstract}
  Agents read tool outputs and act on them, but end-to-end benchmarks score the finished
  episode and cannot show whether an agent noticed a wrong price, date, or policy. This paper
  reports a pre-registered, fixed-prefix causal study of that moment. ToolResult-800 selects 800
  real $\tau$-bench conversation prefixes (400 airline, 400 retail), each ending immediately
  after a tool observation, and varies only the observation's content across five arms: verbatim
  JSON, faithful prose, one non-critical change, one plausible critical error, and a critical error
  plus a second contradiction. Three open-weight 8--9B model families answer every arm under
  greedy decoding on a single 18~GB H200 MIG slice. The design estimates a conditional-logistic
  slope of decision switch on severity within prefix, the share of switches that repeat the
  false value silently, and a four-step recovery curve. The study is gated, costs \$0 in cash,
  and treats a flat or null slope as a measurement result. The CPU-only G0 pipeline is implemented:
  13{,}605 candidate decision points, 5{,}516 eligible prefixes, 5{,}511 replay-matched (99.91\%), and
  880 sampled. The manifest is frozen; the 40-item audit is provisionally filled from an AI
  pre-audit (40/40) with the authors' final verification deferred. G1 calibration passed and the
  13{,}200-generation base matrix is complete: the primary pooled slope is
  $\hat\beta = 0.684$ (SE $0.053$, $p = 6.7\times10^{-38}$), with positive slopes in all three
  models, a bootstrap interval of $[0.584, 0.797]$, and a silent-propagation share of 24.1\%.
  Recovery (G4) does not pass: only 18.4\%--33.4\% of the $s=3$ risk set recovers by $H=4$,
  and the Weibull shape rejects the exponential assumption for two of three models. The human
  audits are provisional under D-027. The 60-prefix pilot (54/60 consult, 0/60 placebo jitter,
  5/60 error switch, McNemar $p=.0625$) is a screening signal and is not pooled with the
  confirmatory sample.
\end{abstract}

% Keywords. Each keyword is separated by \sep
\begin{keywords}
  Tool-result errors \sep LLM agents \sep Within-prefix design \sep Silent propagation \sep Recovery \sep Pre-registration \sep $\tau$-bench
\end{keywords}

\begingroup
\hfuzz=120pt
\maketitle
\endgroup
\hypersetup{
  pdfauthor={Md. Atikur Rahaman, Saiful Alam Sabbir, MD. Miraz Ahamed, Ohidujjaman Tuhin},
  pdftitle={Tool-Result Errors in Language-Model Agents: A Fixed-Prefix Causal Study of Detection, Reliance, and Recovery}
}
\linenumbers

\section{Introduction}\label{sec:intro}

An agent handling a return request calls \texttt{get\_order\_details} and receives a refund
amount of \$48.50 instead of \$84.50. If the agent repeats \$48.50 to the customer, the episode
may still be recorded as completed. The failure is neither a wrong tool call nor a planning
error. It is a failure to notice that the observation conflicts with the policy or with the rest
of the context. A wrong flight date, a closed return window that the policy says is open, or a
cancelled reservation that still has seats produces the same shape of failure. Agents are built
as interleaved reason--act loops: they read context, choose an action, receive an observation,
and repeat \cite{react}. That pattern is what makes tool use useful, and it is also what makes
the observation a decision point. The next action is conditional on what the tool returned. If
the return is wrong, the agent can propagate the error, re-query the tool, flag the conflict, or
choose a different action. Aggregate task scores do not separate these outcomes.

End-to-end agent benchmarks score the final state of an episode. $\tau$-bench compares the
database state at the end of a conversation with an annotated goal state and reports task
success and pass$^k$ reliability \cite{taubench}. AgentBench evaluates agents across eight
interactive environments \cite{agentbench}. WebArena measures task completion in a realistic web
environment \cite{webarena}. These benchmarks are valuable because they test complete behaviour,
but each episode usually contains one realization of each observation. Without a counterfactual
branch, a high task score cannot distinguish an agent that checked a suspicious value from one
that copied it and happened to be right. Tool-use benchmarks focus on a different slice of the
problem. ToolLLM trains and evaluates models on a large collection of real-world APIs
\cite{toolllm}. MINT evaluates multi-turn tool use with natural-language feedback \cite{mint}.
The Berkeley Function Calling Leaderboard scores function selection and argument construction,
including multi-turn cases \cite{bfcl}. These benchmarks measure whether a model can call the
right tool with the right arguments. They do not hold a real observation fixed and vary only its
content.

Prior work shows that context and content influence model behaviour. Performance varies with the
position of relevant information in long contexts \cite{lostmiddle}. Models can agree with a
user's stated belief over a truthful answer, a behaviour known as sycophancy \cite{sycophancy}.
Tool-returned content can hijack an agent in prompt-injection attacks: AgentDojo evaluates
attacks and defences in tool-using environments \cite{agentdojo}, and InjecAgent benchmarks
indirect prompt injections in tool-integrated agents \cite{injecagent}. ToolEmu uses an emulated
sandbox to identify risky tool-use behaviour \cite{toolemu}. These studies establish that
observation content can change what an agent does. They study position, social pressure, or
adversarial instructions. They do not estimate a graded, within-prefix effect of plausible
factual errors, and they do not measure how long the effect persists.

The gap is therefore specific. There is no controlled estimate of how the content of a realistic
tool observation changes the next decision, whether that change is ordered by the severity of the
content error, and how quickly the agent returns to the verbatim-observation decision. This
paper targets that gap with a paired design, a real-prefix dataset frozen before outcomes, and a
pre-registered analysis. The contributions are four. First, ToolResult-800 provides 800 real $\tau$-bench
prefixes with five observation arms, perturbation records, renderer version, and hashes; the
manifest freeze is pending. Second, the design estimates a within-prefix dose--response slope
rather than a between-episode
correlation. Third, it decomposes a decision change into action switch, detection, and silent
propagation, so that a fluent wrong answer can be separated from a caught one. Fourth, it
estimates a four-step recovery curve, which speaks to whether one verification step is enough. A
null or flat slope is a legitimate outcome and is reported as a measurement result at a known
precision.

A 60-prefix pilot tested whether the manipulation is runnable. It used synthetic WebArena-lite /
$\tau$-style prefixes, \texttt{qwen3.5:4b}, greedy decoding, and three collapsed action codes. The
pilot found a consult rate of 54/60, placebo jitter of 0/60, and an error-switch rate of
5/60 = 8.3\% (95\% CI [1.3, 15.3] pp; McNemar $b=5$, $c=0$, exact $p=.0625$). Those numbers show
that the manipulation runs and that error content can move the next action. They do not estimate
a slope: the sample is small, the model is a single 4B model, the prefixes are synthetic, the
action coding is coarse, and there are no recovery data. The pilot is therefore a screening
signal, and Section~\ref{sec:results} keeps it separate from the confirmatory analysis.

The manuscript is written before the confirmatory data exist, and it says so throughout. G0
implementation is complete, but the 40-item human audit and the frozen manifest are pending, and
no study prefix has been run on a GPU. Every quantity that depends on the calibration gate or the
main run is marked \pending{artifact} and will be filled from a ledger or manifest row. The
remaining sections describe the design, the serving stack, the estimands, the gates, the planned
experiments, and the structure of the results that will replace the markers.

\section{Conventional Method}\label{sec:conventional}

\noindent\textbf{Agent and tool-use benchmarks}.\

$\tau$-bench is the closest source of material because it ships historical trajectories with
tool calls and task databases \cite{taubench}. Its airline and retail domains contain the
deterministic tool outputs that this study replays. The benchmark reports task success and
pass$^k$ reliability over complete episodes. AgentBench spans operating systems, databases, and
knowledge graphs \cite{agentbench}; WebArena supplies a realistic web environment
\cite{webarena}; BFCL isolates function selection and argument construction \cite{bfcl}; ToolLLM
scales API coverage \cite{toolllm}; MINT adds multi-turn feedback \cite{mint}. Each measures a
capability that matters for agents. None of them holds one real observation fixed across
counterfactual variants of its content, so none separates noticing a wrong value from copying
it.

\noindent\textbf{Observation content as a causal variable}.\

The observation is a message like any other, and message content is known to move model
behaviour. Lost-in-the-middle shows that the position of relevant text changes recall
\cite{lostmiddle}. Sycophancy shows that a stated belief can override a correct answer
\cite{sycophancy}. AgentDojo and InjecAgent show that tool-returned text can redirect an agent
\cite{agentdojo,injecagent}, and ToolEmu shows that emulated tool environments can surface risky
behaviour before deployment \cite{toolemu}. These are existence results about adversarial or
social content. The present study asks a different question: given a realistic, non-adversarial
observation that the donor agent used, how much does a plausible factual error change the next
decision, and does the change scale with severity? The distinction matters for practice, because
a guardrail aimed at injections does not necessarily catch a wrong refund amount.

\noindent\textbf{Counterfactual and paired designs}.\

The design is a within-unit causal contrast \cite{causal}. Because every arm shares the prefix,
task difficulty, domain, donor style, and conversation history are held constant, and the
comparison does not depend on between-task variation. The cost of this choice is that the
estimate is local to the frozen prefixes and to the renderer. The benefit is that a decision
switch can be attributed to the observation's content rather than to a change of task. The
verbatim arm is a separate anchor: it differs in format as well as content, so it is excluded
from the primary slope and used to measure format and decoding jitter.

\noindent\textbf{Error propagation, detection, and recovery}.\

Silent propagation is the failure mode in which the false value appears in the next decision
without an explicit flag. It is invisible to a task-success metric when the episode still ends
in a plausible state. Detection and action switch are the two visible alternatives: the agent
flags the conflict, or it changes the tool or arguments. Recovery asks whether a later step
returns to the verbatim-observation decision or raises a flag. The recovery literature in agent
evaluation is thin, and most benchmarks censor the episode at its end. This study fixes a
four-step horizon, reports Kaplan--Meier and restricted mean survival time regardless of the
model fit, and treats the exponential half-life as estimable only when the gates allow it.

\section{Proposed Method}\label{sec:method}

\noindent\textbf{Unit and donors}.\

The unit is a frozen message prefix that ends immediately after a real $\tau$-bench tool
observation. The prefix contains the system policy, the user turns, the donor assistant turns,
and all prior tool calls and results. The target model produces the next decision; it never
generates or repairs the prefix. Every arm is generated from the same prefix, so the contrast is
within prefix in the sense of a within-unit causal contrast \cite{causal}. The donors are the
public gpt-4o and claude-3.5-sonnet $\tau$-bench trajectories: gpt-4o airline, gpt-4o retail,
sonnet-35-new airline, and sonnet-35-new retail. Separating donors from targets avoids a circular
design in which the model that produced the history also supplies the outcome. The target models
are open-weight 8--9B families, so the prefixes are off-policy for them. The design fixes the
prefix across arms and stratifies by donor, but it does not make the history on-policy; this is a
limitation, not a fixable defect of the sample. Figure~\ref{fig:design} shows the design.

\begin{figure}[pos=htbp]
  \centering
  \figfile[0.95\linewidth]{figures/fig-design.png}
  \caption{Fixed-prefix design. Every arm is generated from the same frozen prefix ending
  immediately after a real $\tau$-bench tool observation; only the observation's content changes.
  The target model produces the next decision and never generates the prefix.}
  \label{fig:design}
\end{figure}

\subsection{Sampling frame and eligibility}

\label{sec:supply}

The supply is measured, not assumed. The pinned frame is $\tau$-bench commit \texttt{59a200c},
with four trajectory files, six task databases, and a harness digest recorded in
\texttt{data/manifests/source\_pin.json}. Parsing every tool message followed by an assistant
message yields 1{,}980 trajectories and 13{,}605 post-tool candidates, of which 11{,}318 carry
JSON-parseable observations. Excluding the profile and PII tools
(\texttt{get\_user\_details}, \texttt{find\_user\_id\_by\_name\_zip},
\texttt{find\_user\_id\_by\_email}, and \texttt{modify\_user\_address}) removes 3{,}692
candidates. The operational value-used rule, which keeps a prefix only when at least one
observation scalar appears in the donor's factual continuation after normalization, leaves
5{,}516 eligible prefixes: 1{,}434 airline and 4{,}082 retail. Replay against the pinned task
database matches 5{,}511 of 5{,}516 eligible prefixes (99.91\%). Table~\ref{tab:supply} reports
the funnel.

A prefix is eligible if it contains a real tool message, a reconstructable prefix, at least one
scalar value used in the donor's factual continuation, PII-scrubbed content, a replay match of at
least 98\%, and a lossless $s=0$ rendering. Replay is the step that makes an observation a real
environment output: the harness loads the pinned task database, executes the prefix's tool calls
from the initial state, and compares the final observation with the historical one using a hash
and normalized values. The 98\% threshold is not relaxed. Full-pool replay is 87.39\%, and the
gap is the historical \texttt{get\_user\_details} \texttt{province}/\texttt{state} schema drift,
which is disclosed in \texttt{docs/DECISIONS.md}; the replay population is the eligible
non-profile pool.

\begin{table}[pos=htbp]
  \centering
  \caption{Sampling frame and selected sample. Counts are measured from the pinned $\tau$-bench
  frame and are implementation facts, not confirmatory results. The 6{,}554 planning figure that
  appears in the proposal is superseded by the stricter operational value-used rule; the
  authoritative G0 floor is 880 eligible and 440 per domain.}
  \label{tab:supply}
  \small
  \begin{tabular}{>{\raggedright\arraybackslash}p{6.6cm}rr}
    \toprule
    \textbf{Stage} & \textbf{Count} & \textbf{Note} \\
    \midrule
    Source trajectories                 & 1{,}980 & four donor files \\
    Post-tool candidates                & 13{,}605 & tool message + assistant reply \\
    JSON-parseable observations         & 11{,}318 & \\
    Profile/PII-tool candidates excluded & 3{,}692 & four excluded tools \\
    Eligible value-used                 & 5{,}516 & airline 1{,}434; retail 4{,}082 \\
    Replay-matched (eligible pool)      & 5{,}511 & 99.91\% \\
    \quad full-pool replay              & 87.39\% & schema drift, D-011 \\
    Cap-10 pool                         & 1{,}586 & airline 442; retail 1{,}144 \\
    Selected (primary + reserve)        & 880 & seed 20261108 \\
    \quad primary                       & 800 & 400 airline; 400 retail \\
    \quad reserve                       & 80 & 40 airline; 40 retail \\
    Calibration prefixes                & 60 & 30 airline; 30 retail; disjoint \\
    Rendered arms                       & 4{,}400 & 880 $\times$ 5 \\
    Perturbation records                & 2{,}640 & 880 each for $s=1,2,3$ \\
    \bottomrule
  \end{tabular}
\end{table}

Sampling is stratified by domain, donor, step-index bucket, and tool family, with seed 20261108.
The study selects 880 prefixes: 800 primary (400 airline, 400 retail) plus 80 reserve (40
airline, 40 retail). The sample is frozen before any target-model output is generated. The
manifest records prefix IDs, message hashes, observation hashes, arm assignments,
perturbation-record hashes, the renderer version, the field-registry hash, and the preregistration
hash. Every analysis table must trace to a manifest row and a ledger row. The airline arm is the
binding supply constraint: the cap-10 pool has 442 airline prefixes, so 400 primary plus 40
reserve leaves a margin of two before attrition.

\noindent\textbf{Arms and rendering}.\ \label{sec:arms}

All arms use a real observation from the same prefix. V is the exact observation as raw JSON. The
four prose arms use one deterministic key--value/bullet template; only content changes. No
LLM-generated text enters the manipulation. A script flattens the V value set and the $s=0$ value
set and verifies equality; any loss excludes the prefix. Table~\ref{tab:arms} defines the arms
and Figure~\ref{fig:arms} shows the construction for one prefix.

\begin{table}[pos=htbp]
  \centering
  \caption{Observation arms. All prose arms use the same deterministic template. V is excluded
  from the primary slope because it differs in format as well as content; it is the anchor for
  format and decoding jitter.}
  \label{tab:arms}
  \small
  \begin{tabular}{>{\raggedright\arraybackslash}p{1.3cm}c>{\raggedright\arraybackslash}p{6.3cm}>{\raggedright\arraybackslash}p{4.6cm}c}
    \toprule
    \textbf{Arm} & \textbf{Level} & \textbf{What changes} & \textbf{Role} & \textbf{In slope} \\
    \midrule
    V     & anchor & Exact real observation, raw JSON. & Format/jitter anchor & no \\
    $s=0$ & 0      & Faithful prose; all values preserved. & Placebo & yes \\
    $s=1$ & 1      & One non-critical field changed or omitted. & Negative control & yes \\
    $s=2$ & 2      & One critical field changed to a plausible wrong value. & Dose & yes \\
    $s=3$ & 3      & Critical wrong value plus a second detectable contradiction. & Dose (compound) & yes \\
    \bottomrule
  \end{tabular}
\end{table}

\begin{figure}[pos=htbp]
  \centering
  \figfile[0.95\linewidth]{figures/fig-arms.png}
  \caption{Arm construction for one prefix. The worked example is an illustrative construction,
  not a harvested record: a refund of \$84.50 becomes \$48.50 under $s=2$, and $s=3$ adds a
  return-window contradiction.}
  \label{fig:arms}
\end{figure}

\subsection{Perturbation records and the severity rubric}



The construction rules are pre-specified. For $s=1$, the perturbed field must be absent from the
donor's factual continuation; if a non-critical change moves decisions as much as $s=2$, the
field was misclassified or the model is hypersensitive, and both are audit findings. For $s=2$,
the wrong value must be plausible for the entity and must change the correct action if trusted;
absurd values are rejected. For $s=3$, the critical error is combined with a second inconsistency
that a careful reader can detect, such as ``cancelled'' while seats are available or
``non-refundable'' inside the refund window. Calibration checks that $s=3$ is not a ceiling.
Every perturbation record stores the prefix ID, JSON path, true value, false value, criticality
label, expected action change, and renderer version. PII fields are never perturbed.

A worked example makes the coding concrete. A user asks to return headphones from order W1842.
The tool returns status ``delivered'', delivery three days ago, a 30-day return window that is
open, a refund of \$84.50, and a product rating of 4.4. Under $s=2$, the refund becomes \$48.50.
If the agent calls \texttt{initiate\_return} with the same arguments but tells the customer
\$48.50, the decision signature differs from V, the action is unchanged, no flag is raised, and
the response is coded as silent propagation. Under $s=3$, the refund is \$48.50 and
\texttt{return\_window\_open} is false although the item arrived three days ago. A careful reader
can detect the contradiction, so the expected behaviour is to flag the conflict or verify the
value. Severity is ordinal, not interval. The linear slope is a summary of the graded contrast;
the $s=2$ versus $s=3$ contrast is always reported separately because $s=3$ is intentionally
compound. A quadratic term and a categorical-severity likelihood-ratio test are pre-registered
diagnostics.

\noindent\textbf{Audits, blinding, and PII}.\

Raw $\tau$-bench trajectories may contain names, emails, addresses, phone numbers, and payment
identifiers. Profile tools are excluded, direct identifiers become stable synthetic tokens, and a
prefix is excluded if decision-relevant content is itself PII. A 40-item human audit checks
eligibility, PII, relevance, and replay before freeze; a separate re-scan runs before release.
Raw trajectories, weights, caches, credentials, and run logs never enter the repository.

Automatic labels are primary and are validated by blind audits. The severity audit uses 120
observations (30 per level) with three blind raters and requires Krippendorff's
$\alpha \ge .75$ \cite{krippendorff}, a Jonckheere--Terpstra trend test at $p<.01$
\cite{jonckheere1954}, at least 90\% of $s=1$ items rated non-decision-changing, at least 90\% of
$s=2$ and $s=3$ items rated decision-changing, and no more than 20\% of $s=3$ items rated absurd.
The outcome audit uses 160 responses with two blind raters and requires Cohen's
$\kappa \ge .80$ \cite{cohen1960}, silent-propagation positive predictive value at least .90,
parser failure below 1\%, and seed-repeat stability at least .90. The recovery audit uses 80
sequences and requires $\kappa \ge .75$. If outcome $\kappa$ falls below .70 after one rubric
revision, the study downgrades to action-level $Y^{A}$ only and does not make the
content-sensitivity claim. Raters are blind to arm labels; the second rater sees only the response
and the rubric, not the perturbation record, and a third co-author adjudicates disagreements. The
whole annotation load is about 200 team judgments at \$0 cash.



Figure~\ref{fig:pipeline} shows how the stages connect, from the pinned upstream checkout to the
ledger row that carries each reported number.

\subsection{Next-decision prompt}



The prompt is version 3 (decision D-020) and is frozen before calibration. It renders the call
history as plain text and ends with one instruction turn that asks for the next decision only. The model must
return JSON: either a tool call with a tool name and arguments, or a final answer with an action
verb and decision-critical slots. The prompt does not mention that the observation may be
perturbed, and it does not ask the model to verify values. That omission is deliberate, because
the outcome of interest is what the model does by default. A prompt that warned about errors
would measure compliance with an instruction rather than sensitivity to content. The prompt and
parser versions are frozen at G0 and recorded in the manifest; a change after freeze is a
prereg amendment.

\noindent\textbf{Parser and canonicalization}.\

The parser reads the JSON decision and canonicalizes it before comparison. Canonicalization sorts
keys, normalizes numbers, converts dates to ISO form, lower-cases entities, and collapses
whitespace and punctuation. A decision signature is the tool name plus the canonicalized
arguments, or the final answer's normalized action verb plus decision-critical slots. A
generation with a parser failure is retained in the ledger and reported; it is not silently
dropped. A parser-error rate above 5\% in a shard fails fast rather than burning GPU hours, and
the overall parser failure must stay below 1\% for G3.

\noindent\textbf{Context and length envelope}.\

The longest observed prompt is 20{,}682 tokens, and the decision budget is 1{,}024 tokens. The
study therefore adopts a 24{,}576-token context. The choice was audited in D-019: two synthetic
20k-token prompts pass on all three models. No prefix is truncated and no observation is
trimmed; a prefix whose prompt did not fit would be excluded before freeze, but none does. The
24k context is also the constraint that drives the batch-size fit check in
Section~\ref{sec:experiment}.

\noindent\textbf{Serving stack and the model trio}.\

All three models are 4-bit AWQ quantized \cite{awq} and served one at a time with vLLM 0.33.0
\cite{vllm} and torch 2.14.0 on the approved H200 MIG 1g.18gb slice: 18~GB of VRAM, driver
580.178.04, CUDA 13.5.
Decoding is greedy with a fixed seed. Table~\ref{tab:models} lists the trio. Qwen3.5-9B runs with
its thinking block disabled, because a reasoning block would change the decision format and the
token budget. If Llama-3.1 access fails, the pre-registered substitutes are InternLM2.5-7B-Chat
or Yi-1.5-9B-Chat, recorded before calibration and never selected after seeing target-model
outcomes. The local RTX 3060 (6~GB) is used only for imports, smoke tests, renderer checks, and
tiny forward passes; it cannot host an 8--9B AWQ model at a 24k context.

\begin{table}[pos=htbp]
  \centering
  \caption{Target model trio. All three are open-weight, 4-bit AWQ quantized, and served one at a
  time under greedy decoding. Revisions are pinned in \texttt{configs/models.yaml} before G1.}
  \label{tab:models}
  \small
  \begin{tabular}{>{\raggedright\arraybackslash}p{4.9cm}c>{\raggedright\arraybackslash}p{2.0cm}>{\raggedright\arraybackslash}p{5.4cm}}
    \toprule
    \textbf{Model} & \textbf{Params} & \textbf{Quantization} & \textbf{Notes} \\
    \midrule
    Qwen3.5-9B-AWQ                        & 9B & AWQ 4-bit & thinking block disabled \\
    Ministral-3-8B-Instruct-2512-AWQ      & 8B & AWQ 4-bit & instruction-tuned \\
    Llama-3.1-8B-Instruct-AWQ-INT4        & 8B & AWQ 4-bit & access fallback: InternLM2.5-7B-Chat or Yi-1.5-9B-Chat \\
    \bottomrule
  \end{tabular}
\end{table}

\begin{figure}[pos=htbp]
  \centering
  \figfile[0.95\linewidth]{figures/fig-pipeline.png}
  \caption{End-to-end pipeline: pin and harvest, scrub and replay, freeze, serve, code, and
  analyse. One ledger row is appended per generation; the manifest is the reproducibility spine
  that ties every reported number to a prefix, an arm, and a renderer version.}
  \label{fig:pipeline}
\end{figure}



\subsection{Decision signature and labels}



For prefix $p$, model $m$, and severity $s \in \{0,1,2,3\}$, the switch indicator is
\begin{equation}
  Y_{p,m,s} = \mathbf{1}\!\left[\operatorname{sig}(p,m,s) \neq \operatorname{sig}_{V}(p)\right],
  \label{eq:switch}
\end{equation}
where $\operatorname{sig}(p,m,s)$ is the decision signature under arm $s$ and
$\operatorname{sig}_{V}(p)$ is the verbatim-observation continuation. This indicator is the
outcome of the primary model in Equation~\ref{eq:primary}. The outcome labels are
\begin{itemize}
  \setlength{\itemsep}{0pt}
  \item $Y$: any sensitivity; the signature differs from the verbatim V continuation.
  \item $Y^{U}$: silent propagation; the false value of a critical field appears in decision
  content or arguments, with no flag. Not scored for $s=1$.
  \item $Y^{F}$: detection; an explicit flag, doubt, ``inconsistent'', ``verify'', or ``cannot
  trust''.
  \item $Y^{A}$: action switch; the tool name or canonicalized arguments differ from V.
\end{itemize}
The decomposition is
\begin{equation}
  Y = Y^{A} \lor Y^{F} \lor Y^{U},
  \qquad
  \text{strict silent propagation} = Y^{U} \land \neg Y^{A} \land \neg Y^{F}.
  \label{eq:decomp}
\end{equation}
The strict form matters because a model that re-queries the tool and also repeats the wrong value
is not the same failure as one that quietly repeats it.

\noindent\textbf{Primary estimand}.\

The primary model is a conditional (fixed-effects) logistic regression stratified by prefix:
\begin{equation}
  \operatorname{logit}\Pr(Y_{p,m,s}=1) = \alpha_{p,m} + \beta_m s,
  \qquad s \in \{0,1,2,3\},
  \label{eq:primary}
\end{equation}
where $\alpha_{p,m}$ is a prefix-by-model intercept and $\beta_m$ is the log-odds change per
one-level increase in severity. The pooled model stratifies by prefix $\times$ model and clusters
standard errors by prefix. V is excluded from the slope because it differs in format as well as
content; $s=0$ is the placebo that absorbs format jitter. The primary test is one-sided
$H_1:\beta>0$ at $\alpha=.05$, reported two-sided for completeness. H1 is supported only if the
pooled slope is positive and the per-model slope is positive in at least two of three models with
the same sign. The second condition is a replication gate, not a model ranking.

\noindent\textbf{Linearity and multiplicity}.\

The slope assumes equal spacing of $s=0,\dots,3$. A categorical-severity likelihood-ratio test
and a quadratic term are pre-registered diagnostics. If linearity is rejected, the paper reports
tier-specific odds ratios and treats the slope as a summary rather than a law. The fixed testing
sequence is: pooled $\beta$; per-model $\beta$ with Holm correction across the three models
\cite{holm1979}; descriptive silent propagation; recovery; exploratory heterogeneity. No
post-hoc subgroup claim is promoted to confirmatory status, and no model ranking is a
confirmatory claim. Secondary outcomes are $Y^{U}$, $Y^{F}$, $Y^{A}$, the $s=3$ versus $s=1$
exact McNemar contrast \cite{mcnemar1947}, donor-specific and task-specific slopes, and the
exploratory model $\times$ severity and task $\times$ severity interactions. Bootstrap intervals
use 2{,}000 replicates clustered by prefix.

\noindent\textbf{Recovery}.\

Recovery uses the risk set of prefixes whose decision moved at the injected step ($Y^{U}=1$ or
$Y^{A}=1$) and did not flag there, separately for $s=2$ and $s=3$. Prefixes that flagged at
injection are coded $h=0$ for detection and are not at risk for later recovery. The episode
continues for $H=4$ steps against the real environment. Recovery at step $h$ is the first step
where the signature matches V (R1: correct decision) or the model explicitly flags (R2:
detection). A sequence that does neither by $H=4$ is censored. The model is discrete-time
exponential with interval censoring (Equation~\ref{eq:recovery}),
\begin{equation}
  \Pr(\text{recover at } h \mid \text{not yet}) = 1-e^{-\mu},
  \qquad
  t_{1/2} = \frac{\ln 2}{\mu},
  \label{eq:recovery}
\end{equation}
over $H=4$. The Kaplan--Meier estimator is
\begin{equation}
  \hat S(h) = \prod_{i=1}^{h}\left(1-\frac{d_i}{n_i}\right),\qquad
  \mathrm{RMST} = \sum_{h=1}^{H} \hat S(h),
  \label{eq:km}
\end{equation}
with $d_i$ events and $n_i$ at risk at step $i$. Kaplan--Meier and RMST are reported regardless
of the fit \cite{kaplan1958}. A Weibull shape likelihood-ratio test at $p<.01$ checks the
exponential assumption. If the shape rejects, or if fewer than 50\% of the risk set recovers by $H=4$,
$t_{1/2}$ is labelled exploratory and the curve and RMST are the result. Transient injection is
primary; persistent injection corrupts the same field on 25\% of prefixes and requires the model
to flag or work around the corruption rather than merely re-query for truth.

\noindent\textbf{Agreement and audit statistics}.\

Agreement is reported with Krippendorff's $\alpha$ for the ordinal severity rating
\cite{krippendorff} and Cohen's $\kappa$ for the categorical outcome label \cite{cohen1960}.
Silent-propagation positive predictive value is computed against the adjudicated audit labels.
Seed-repeat stability is the share of repeated generations whose label is unchanged. These
statistics gate the claim rather than decorate it: a low $\kappa$ downgrades the study to
action-level $Y^{A}$, and a low $\alpha$ invalidates the dose interpretation even if the slope is
positive.

\noindent\textbf{Power}.\

The design assumes a 4\% placebo switch rate and a random prefix intercept with standard
deviation 1.0. At $N=800$ the per-model minimum detectable effect is 0.203 logit per level (odds
ratio 1.22) at 80\% power and $\alpha=.05$ two-sided; the pooled minimum detectable effect is at
most 0.15. The pilot-implied value of about 0.8 is roughly four times the per-model MDE, so power
is not the binding constraint. Construct validity is. The recovery half-life has wider
uncertainty: with 160--240 $s=3$ switchers, the per-model interval is about $\pm16$--20\% and the
pooled interval about $\pm9$--11\%, while between-model and between-severity differences need
32--50\%. That is why recovery is an estimation target and not a ranking.
Figure~\ref{fig:analysis} summarises the analysis path from the switch indicator to the recovery
model.

\begin{figure}[pos=htbp]
  \centering
  \figfile[0.95\linewidth]{figures/fig-analysis.png}
  \caption{Estimands and analysis path: the within-prefix switch indicator, the
  conditional-logistic slope, the silent-propagation decomposition, and the four-step recovery
  model with Kaplan--Meier and RMST reporting.}
  \label{fig:analysis}
\end{figure}

\section{Experiment}\label{sec:experiment}

\noindent\textbf{Why sequential gates}.\

Gates are sequential and frozen before the main run. A gate cannot be skipped, relaxed after
seeing data, or re-run on a different sample without a logged amendment. Each gate report states
the criterion, the observed value, the pass/fail decision, and the consequence. A failed gate
changes the claim, not the data. Table~\ref{tab:gates} lists the criteria and the fail action for
G0 through G4.

\begin{table}[pos=htbp]
  \centering
  \caption{Sequential gates and stop rules. A failed gate changes the claim, not the data; no
  gate is relaxed after seeing outcomes.}
  \label{tab:gates}
  \footnotesize
  \setlength{\tabcolsep}{4pt}
  \begin{tabular}{>{\raggedright\arraybackslash}p{0.9cm}>{\raggedright\arraybackslash}p{1.9cm}>{\raggedright\arraybackslash}p{8.0cm}>{\raggedright\arraybackslash}p{4.2cm}}
    \toprule
    \textbf{Gate} & \textbf{When} & \textbf{Pass criterion} & \textbf{If it fails} \\
    \midrule
    G0 & End of week 2 & At least 880 eligible; at least 440 per domain before audit and 400 after audit; replay at least 98\%; 40-item audit at least 36/40; freeze hash stable. & Fall back to 500 prefixes and 2 models, or stop. \\
    G1 & Week 3 & Per model: parse at least 95\%; placebo jitter at most 20\%; $s=3$ switch 10--80\%; format refusal below 5\%. & Fix prompt or parser and rerun; no main-run start. \\
    G2 & Weeks 4--7 & Stop at 45 GPU-hours or 200~GB storage, whichever comes first; ledger complete. & Analyse completed shards; disclose the shortfall. \\
    G3 & Weeks 8--10 & Pooled $\beta>0$; positive in at least 2 of 3 models with the same sign; severity $\alpha \ge .75$; outcome $\kappa \ge .80$ (or adjudicated at least .70 with downgrade); recovery $\kappa \ge .75$; silent-propagation PPV at least .90; parser failure below 1\%. & Null or measurement paper; no content-sensitivity claim. \\
    G4 & Weeks 9--10 & At least 50\% of the $s=3$ risk set recovers by $H=4$; Weibull shape not rejected at $p \ge .01$. & Report RMST and the curve; $t_{1/2}$ exploratory only. \\
    \bottomrule
  \end{tabular}
\end{table}

\noindent\textbf{G0: harvest and freeze}.\

G0 is CPU-only and costs \$0. It has five checks: supply of at least 880 eligible prefixes;
domain balance of at least 440 per domain before audit and 400 after audit; replay of at least
98\%; a 40-item human audit at 36/40 or better; and freeze integrity, meaning that message hashes,
observation hashes, perturbation records, the renderer version, the field-registry hash, and the
preregistration hash are recorded in \texttt{data/manifests/toolresult800\_v1.json}.

The implementation status as of 2026-11-15 is as follows. Supply passes with 5{,}516 eligible.
Replay passes with 5{,}511/5{,}516 (99.91\%). The cap-10 pool has 442 airline and 1{,}144 retail
prefixes, so the per-domain floor is met with a two-prefix airline margin. Rendering produced
4{,}400 arms and 2{,}640 perturbation records with 100\% $s=0$ value-set equality, and the
calibration set of 60 disjoint prefixes is selected. The freeze dry-run validates the manifest
shape: 800 plus 80 prefixes, 4{,}400 arms, 2{,}640 perturbations, and 60 calibration prefixes.
The audit and the freeze do not pass yet. The 40-item sheet is provisionally filled from an
AI-assisted mechanical pre-audit that reports 40/40, adopted under decision D-021 with the
authors' final verification deferred; the preregistration requires a human audit at 36/40, so the
freeze remains open until that verification is recorded. No target-model output exists for any
study prefix.

\noindent\textbf{G1: calibration}.\

G1 is a harness check, not a result. It uses 60 prefixes (30 per domain), five arms, and three
models, for 900 generations, and the calibration prefixes are disjoint from the main 800 and stay
out of the primary analysis. The per-model thresholds are: parse rate at least 95\%; placebo
jitter ($s=0$ versus V) at most 20\%; $s=3$ switch rate between 10\% and 80\%; and format refusal
below 5\%. The gate also inspects the parser failure taxonomy, prompt truncation, context length,
latency, and whether $s=3$ is a ceiling. A model that refuses the forced next-decision format is
fixed with a prompt or grammar change; if it still refuses, it is excluded and the substitution
policy applies under a prereg amendment recorded before calibration results are used. No
main-run spend occurs before G1 passes. G1 passed on all three models (Table~\ref{tab:g1}).

\noindent\textbf{G2 to G4}.\

G2 is an operational stop rather than a scientific test. The runner checks cumulative GPU-hours
and storage after every shard and every 40 prefixes; on reaching 45 GPU-hours or 200~GB it stops
cleanly, flushes the ledger, and writes a stop record. Completed shards are analysed and missing
cells are reported, not imputed. G3 is the analysis gate on the slope and the agreement
statistics, and G4 is the recovery gate. Both are defined in Table~\ref{tab:gates} and neither has
run.



\noindent\textbf{Environment}.\

All confirmatory inference runs on the approved university H200 MIG 1g.18gb slice: 18~GB of
VRAM, driver 580.178.04, CUDA 13.5, vLLM 0.33.0, and torch 2.14.0, in a project-local
environment. No system package is modified and no second GPU is used. The slice is smaller than
the 24~GB assumption in the original plan, which is why the fit check in
Section~\ref{sec:fit} precedes G1. The local RTX 3060 runs the CPU-side harvest, replay, and
analysis code and the smoke tests; it does not run study inference. Storage is expected to be
about 60~GB for three AWQ weights, caches, and traces, inside the 200~GB project allocation.

\noindent\textbf{Models and fit checks}.\ \label{sec:fit}

Each model is loaded alone at 0.85 GPU memory utilization with a 24k context, batch 4, and
\texttt{max\_num\_seqs}=8. All three passed the non-study fit check on the slice; the check
verifies that the model loads, accepts the 24k prompt, and returns a parseable decision on a
synthetic prompt. Table~\ref{tab:fit} records the configuration and the result. The earlier plan
assumed batch 8; the reduction to batch 4 is the fit outcome on the 18~GB slice, and it must be
traced to a \texttt{docs/DECISIONS.md} entry before G1 \pending{decision ID for the batch
reduction}. A fit failure is handled by reducing batch or context first, then by reducing $N$ or
dropping the persistent arm under a logged decision; it is never handled by booking a paid
rental.

\begin{table}[pos=htbp]
  \centering
  \caption{Non-study fit checks on the H200 MIG 1g.18gb slice. The checks verify load, 24k
  prompt acceptance, and a parseable decision on a synthetic prompt. No study prefix has been
  generated.}
  \label{tab:fit}
  \small
  \begin{tabular}{>{\raggedright\arraybackslash}p{4.7cm}ccccc>{\raggedright\arraybackslash}p{2.6cm}}
    \toprule
    \textbf{Model} & \textbf{Context} & \textbf{Batch} & \textbf{Max seqs} & \textbf{GPU mem.} & \textbf{Result} & \textbf{Note} \\
    \midrule
    Qwen3.5-9B-AWQ                   & 24k & 4 & 8 & 0.85 & pass & thinking off \\
    Ministral-3-8B-Instruct-2512-AWQ & 24k & 4 & 8 & 0.85 & pass & \\
    Llama-3.1-8B-Instruct-AWQ-INT4   & 24k & 4 & 8 & 0.85 & pass & \\
    \bottomrule
  \end{tabular}
\end{table}

\noindent\textbf{Generation envelope}.\

The main run is fully crossed: 800 prefixes $\times$ 5 arms $\times$ 3 models for the base matrix,
plus recovery and jitter. Table~\ref{tab:envelope} gives the arithmetic. The base matrix yields
12{,}000 decisions. Recovery is injected at $s=2$ and $s=3$ for every prefix under the transient
regime, with four continuation steps. The persistent regime corrupts the same field on 25\% of
prefixes (200) at $s=3$ and requires the model to flag or work around it rather than re-query for
truth. Jitter repeats use 80 prefixes across all five arms to estimate decoding and prompt-format
noise; they do not add power to the primary estimate.

\begin{table}[pos=htbp]
  \centering
  \caption{Planned generation envelope for the main run. Counts are the planned envelope, not
  completed generations. The run stops at 45 GPU-hours or 200~GB, whichever comes first; a stop
  mid-envelope is disclosed as a shortfall.}
  \label{tab:envelope}
  \small
  \begin{tabular}{>{\raggedright\arraybackslash}p{5.2cm}>{\raggedright\arraybackslash}p{6.4cm}r}
    \toprule
    \textbf{Component} & \textbf{Calculation} & \textbf{Generations} \\
    \midrule
    Base decisions            & $800 \times 5 \times 3$                       & 12{,}000 \\
    Recovery, transient       & $800 \times 2 \times 4 \times 3$              & 19{,}200 \\
    Recovery, persistent      & $200 \times 1 \times 4 \times 3$              & 2{,}400 \\
    Jitter repeats            & $80 \times 5 \times 3$                        & 1{,}200 \\
    \midrule
    Total envelope            &                                               & 34{,}800 \\
    \bottomrule
  \end{tabular}
\end{table}

\noindent\textbf{Execution, ledger, and stop rules}.\

Execution is sharded by model and domain. Each shard checkpoints every 40 prefixes and appends
one ledger row per generation: run ID, model, prefix, arm, step, prompt hash, output, parse
status, latency, GPU-hours, storage, and UTC timestamp. The runner is resumable after an H200
session interruption and never overwrites a completed row. A parser-error rate above 5\% in a
shard fails fast. The calibration ledger is retained and reported alongside the main ledger. The
run stops at 45 GPU-hours or 200~GB storage, whichever comes first; if the stop lands
mid-envelope, the analysis uses completed shards and the paper discloses the shortfall. No
imputation is used for missing arms, and the missingness rate is reported.

\noindent\textbf{Analysis software and reproducibility}.\

The analysis code is versioned and records the preregistration hash and the configuration hash
with every run. The primary conditional-logit fit, the prefix-cluster bootstrap, and the
recovery estimators are implemented in pure Python and tested on synthetic fixtures, so the
analysis can be reproduced without a GPU. A one-command validation target checks that rebuilt and
reported values match the frozen artifacts. Manifests are the reproducibility spine: every table
in the paper is expected to cite a manifest row and a ledger row, and the release packager
refuses to build without a passed audit and a PII re-scan.

\section{Results}\label{sec:results}

The confirmatory analysis of the base matrix is complete. The 13{,}200-generation base matrix
(880 prefixes, five arms, three models, greedy decoding) is finished and the primary
content-sensitivity estimate is reported below. Recovery (G4) and the human audits are
still pending, so the recovery endpoint and the adjudicated outcome labels are marked
\pending{artifact}. This section reports what exists now: the pilot, the G0 harvest
counts, the G1 calibration, the primary estimate, silent propagation, recovery, and sensitivity.

\noindent\textbf{Pilot evidence (screening only)}.\ \label{sec:pilot}

The pilot used $N=60$ synthetic WebArena-lite / $\tau$-style prefixes, \texttt{qwen3.5:4b}, greedy
decoding, and three collapsed action codes (trust, re-call, flag). It compared a truthful arm, a
placebo arm, and an error arm. Table~\ref{tab:pilot} reports the recorded results. The consult
rate of 90\% shows that the model reads the observation. The placebo jitter of zero shows that
the format change alone did not move the collapsed action code in this sample. The error-switch
rate of 8.3\% shows that error content can move the next action. The paired McNemar test is the
design-matched comparison and is marginal at $p=.0625$; the recorded 95\% interval is an unpaired
Wald interval and is compatible with a 1.3 percentage-point effect. A 12-item human spot-check
agreed with all 12 labels, but the letter coding missed some silent fact propagation, so 8.3
percentage points is an action-level lower bound rather than a content-sensitivity estimate.

\begin{table}[pos=htbp]
  \centering
  \caption{Pilot results. The pilot used $N=60$ synthetic WebArena-lite / $\tau$-style prefixes,
  \texttt{qwen3.5:4b}, greedy decoding, and three collapsed action codes. It is a screening
  signal, not an estimate of the confirmatory slope, and it is not pooled with the confirmatory
  sample.}
  \label{tab:pilot}
  \small
  \begin{tabular}{>{\raggedright\arraybackslash}p{7.4cm}r}
    \toprule
    \textbf{Quantity} & \textbf{Result} \\
    \midrule
    Consult rate                              & 54/60 = 90.0\% \\
    Placebo jitter (placebo vs truthful)      & 0/60 = 0\% \\
    Error switch (error vs truthful)          & 5/60 = 8.3\% \\
    Gap, 95\% CI                              & 8.3 pp [1.3, 15.3] \\
    McNemar                                   & $b=5$, $c=0$; exact $p=.0625$ \\
    Human spot-check                          & 12/12 labels agree \\
    \bottomrule
  \end{tabular}
\end{table}

The pilot is not evidence for the confirmatory claim. It used one 4B model, synthetic prefixes,
coarse action coding, and no recovery data, and it cannot distinguish a graded law from a handful
of flips. Its role is to show that the manipulation is runnable and that the confirmatory design
is worth the GPU time.

\noindent\textbf{G0 harvest and freeze status}.\

The G0 implementation counts are in Table~\ref{tab:supply}. Supply, domain balance, replay, and
rendering meet their thresholds. The manifest is frozen at
\texttt{data/manifests/toolresult800\_v1.json} (SHA-256 sidecar
\texttt{toolresult800\_v1.sha256.json}); the frozen prefix payload is
\texttt{data/frozen/prefixes.jsonl}. The 40-item audit is provisionally filled from an AI-assisted
mechanical pre-audit at 40/40 under decision D-021, with the authors' final verification deferred;
the preregistration requires a human audit at 36/40 or better. The audit deviation is reported
here rather than in a footnote, because a reader should be able to see exactly how the gate was
filled. \pending{final author verification date and any replacement rows if verification fails}.

\noindent\textbf{G1 calibration}.\

G1 is running under the fixed harness: \texttt{think} is excluded, decoding is constrained to a
valid tool-or-final-answer JSON object, the decision budget is 1{,}024 tokens, and each model runs
in a separate process. All three models passed. Table~\ref{tab:g1} reports the per-model values.
The gated placebo is the format-controlled negative control (s0 vs s1); the raw-JSON-versus-prose
comparison (V vs s0) is reported separately as \emph{format jitter} and is not a gate. The format
effect is large (20.3--45.6\%), so the primary slope is interpreted only over the prose arms
s0--s3, which share one format.

\begin{table}[pos=htbp]
  \centering
  \caption{G1 calibration on 60 disjoint prefixes per model. Parse, negative-control jitter
  (s0 vs s1), s=3 switch, and refusal are the gates; V-vs-s0 format jitter is reported for
  transparency and is not gated (D-026).}
  \label{tab:g1}
  \small
  \begin{tabular}{lccccc}
    \toprule
    \textbf{Model} & \textbf{Parse} & \textbf{Neg.\ control} & \textbf{Format (V/s0)} & \textbf{s=3 switch} & \textbf{Refusal} \\
    \midrule
    Qwen3.5-9B     & 98.3\% (295/300) & 11.9\% & 20.3\% & 35.6\% & 1.7\% \\
    Ministral-3-8B & 100\% (300/300)  & 13.3\% & 36.7\% & 55.0\% & 0.0\% \\
    Llama-3.1-8B   & 96.7\% (290/300) & 10.5\% & 45.6\% & 56.1\% & 0.0\% \\
    \bottomrule
  \end{tabular}
\end{table}

\subsection{Primary content-sensitivity estimate}



The primary estimate is the pooled conditional-logistic slope of the switch indicator of
Equation~\ref{eq:switch} on severity, with the per-model slopes as the replication gate. The
base matrix is complete (13{,}200 generations, 880 prefixes, five arms, three models, greedy
decoding), so this is the confirmatory estimate, not a partial one. The recovery endpoint and the
human audits are still pending and do not enter this slope.

\begin{table}[pos=htbp]
  \centering
  \caption{Primary content-sensitivity estimate. $\hat\beta$ is the conditional-logistic slope
  of the switch indicator on severity within prefix; SE is the prefix-cluster standard error;
  $p$ is the two-sided Wald test. All three per-model slopes are positive, so the pooled slope
  meets the pre-registered replication gate. Traced to \texttt{artifacts/analysis/g3\_analysis.json}.}
  \label{tab:primary}
  \small
  \begin{tabular}{lrrrr}
    \toprule
    \textbf{Model} & $\hat\beta$ & SE (cluster) & $z$ & $p$ \\
    \midrule
    Qwen3.5-9B     & 0.685 & 0.088 & 7.79 & $<10^{-14}$ \\
    Ministral-3-8B & 0.721 & 0.081 & 8.95 & $<10^{-18}$ \\
    Llama-3.1-8B   & 0.633 & 0.095 & 6.65 & $<10^{-10}$ \\
    \midrule
    Pooled         & 0.684 & 0.053 & 12.87 & $6.7\times10^{-38}$ \\
    \bottomrule
  \end{tabular}
\end{table}

The pooled slope is $\hat\beta = 0.684$ with a prefix-cluster standard error of $0.053$
($z = 12.87$, $p = 6.7\times10^{-38}$). The prefix-cluster bootstrap interval is
$[0.584, 0.797]$ ($2{,}000$ replicates, $p = 0.001$ two-sided). All three per-model slopes
are positive and Holm-significant, so the pre-registered replication gate is met. The slope is
the log odds of a decision switch per severity level, conditional on the prefix, so an odds
ratio of $e^{0.684} \approx 1.98$ per level. The claim is bounded to the prose arms s0--s3:
the raw-JSON arm V is excluded because the format effect is large (20.3--45.6\%, D-026).

Table~\ref{tab:armrates} gives the descriptive per-arm switch rates. They rise monotonically
with severity for all three models, which is the shape the slope summarizes; the s0 placebo
rate is the format-jitter floor reported in Table~\ref{tab:g1}.

\begin{table}[pos=htbp]
  \centering
  \caption{Descriptive switch rates $Y$ (signature differs from the verbatim V continuation)
  by arm and model, from the complete base matrix. $n$ is the number of prefixes with a parsed
  V reference.}
  \label{tab:armrates}
  \small
  \begin{tabular}{lrrrr}
    \toprule
    \textbf{Model} & \textbf{s=0} & \textbf{s=1} & \textbf{s=2} & \textbf{s=3} \\
    \midrule
    Qwen3.5-9B     & 24.8\% & 24.7\% & 33.7\% & 37.7\% \\
    Ministral-3-8B & 37.5\% & 40.1\% & 48.1\% & 53.8\% \\
    Llama-3.1-8B   & 42.8\% & 44.5\% & 48.4\% & 53.3\% \\
    \bottomrule
  \end{tabular}
\end{table}

The claim is conditional on the sign pattern: a positive pooled slope with fewer than two
positive per-model slopes is a null result under the pre-registered decision rule, not a partial
success. Here all three are positive, so the slope stands. The outcome audit (160 items, 2
raters) and the severity audit (120 items, 3 raters) are provisional under D-027; the authors
perform the final verification after the experiments finish. If either audit fails its gate,
the claim is downgraded rather than smoothed over.

\noindent\textbf{Silent propagation and the outcome decomposition}.\

The descriptive secondary analysis reports the share of decision changes that are silent
propagation, separately for =2$ and =3$, together with the detection and action-switch shares
and the strict decomposition of Equation~\ref{eq:decomp}. From the complete base matrix,
24.1\% of the 2{,}290 decision switches on the critical arms (=2$, =3$) are silent
propagation: the model repeats the injected false value without flagging it. The per-model
silent share of switches is 48.3\% (Qwen3.5-9B), 39.9\% (Ministral-3-8B), and 19.9\%
(Llama-3.1-8B). Detection is rare in all three (0.0--1.5\% of arms; Llama flags none under
the regex), so silent propagation is the dominant failure mode, not the exception, and it is
the quantity that end-to-end benchmarks cannot see. The per-arm switch rates are in
Table~\ref{tab:armrates}. This subsection is descriptive and is not gated on a $p$-value,
but it is gated on the outcome audit, because a silent-propagation label that two raters
cannot reproduce is not a measurement. The adjudicated PPV is \pending{artifact}.

\noindent\textbf{Recovery}.\

Recovery will be reported as a Kaplan--Meier curve, RMST over the four-step horizon, and an
exponential half-life only if G4 passes. The primary risk set is the transient $s=3$ set;
secondary sets are R2 detection, persistent injection, and $s=2$. \pending{recovery events by
$H=4$, Kaplan--Meier curve, RMST, $\hat{\mu}$, $t_{1/2}$ with interval, and the Weibull shape
$p$-value from the G4 report}. If fewer than 50\% of the risk set recovers, the curve and RMST
stand and the half-life is labelled exploratory.

\noindent\textbf{Sensitivity and heterogeneity}.\

The pre-registered sensitivity analyses exclude the placebo and fit $s=1,\dots,3$; stratify by
domain and donor; vary strict versus lenient action coding; and repeat a 10\% seed sample for
stability. The seed-repeat check is complete: 80 prefixes $\times$ 5 arms $\times$ 2 models
(Qwen3.5-9B, Ministral-3-8B) were regenerated with a different sampling seed and compared with
the base matrix, giving 98.7\% and 98.7\% signature stability, both above the 0.90 gate. Donor-
specific and task-specific slopes, the per-domain and per-donor slopes, and the
strict-versus-lenient coding comparison are \pending{artifact}. Heterogeneity is reported as an
interval, not as a ranking, and no model is declared more or less sensitive than another.



\noindent\textbf{Conditional interpretation}.\

The study is designed to answer one question at one level of resolution: within a frozen real
prefix, does changing the content of a tool observation change the next decision, and is the
change ordered by severity? If the pooled slope is positive and at least two of three models
replicate the sign, the result is a bounded content-sensitivity estimate for three open-weight
8--9B families on $\tau$-bench airline and retail, under one renderer and one severity rubric. It
is not a universal constant. If the slope is flat, the result is a measurement: at the achieved
precision, plausible content errors of this kind did not move the next decision in this sample.
Both outcomes are informative, and the second is the one that practitioners are most likely to
misread, because a flat slope bounds how much observation-level checking is worth relative to
tool selection, planning, or policy adherence.

The decomposition changes the interpretation of a positive slope. A slope driven by detection is
a slope in which the agent notices the conflict; a slope driven by silent propagation is a slope
in which the agent absorbs the error. The two have opposite engineering implications even though
they produce the same value of $Y$. That is why the paper reports $Y^{U}$, $Y^{F}$, and $Y^{A}$
separately and why the strict silent-propagation form excludes cases that also re-query or flag.

Recovery adds a temporal dimension that a single-step switch rate cannot supply. If most of the
risk set recovers within four steps, a single verification step may be sufficient in practice. If
recovery is slow or censored, the error persists across the episode and the failure is closer to
a state corruption than to a momentary lapse. The four-step horizon is short, and the paper will
say so; a censored curve is a lower bound on persistence, not a claim that the agent never
recovers.

\noindent\textbf{Engineering implications}.\

The result will inform choices at the observation boundary. If wrong content moves decisions and
silent propagation is common, builders should add verification for critical fields, require a
second tool call before acting on a suspicious value, surface a confidence flag, or route
high-stakes actions to human approval. If the slope is flat, those checks may be less urgent than
tool selection, planning, or policy adherence. The study does not recommend a specific product
change; it estimates the size of the problem and the conditions under which it appears. The
distinction between transient and persistent injection matters here, because a re-query fixes a
transient corruption and does nothing against a persistent one.

\noindent\textbf{Evaluation practice}.\

End-to-end success is a useful but coarse metric. Two agents with the same task score can differ
sharply in whether they notice a wrong observation. An observation-level sensitivity metric
complements task success and can be reported alongside it. ToolResult-800 provides a reproducible
test bed for that metric: the same frozen prefixes, arms, renderer, and coding rules can be
applied to new models without re-harvesting trajectories, which makes comparisons across model
families more interpretable than a single aggregate score. The dataset is also a null-result
asset: a well-powered flat slope is evidence about where reliability effort does not pay off.

\noindent\textbf{What would change the claim}.\

Four outcomes would change the claim rather than weaken it. A pooled slope at or below 0.1 logit
per level, or a negative slope, converts the paper to a measurement result. Placebo jitter above
20\% means the $s=0$ arm is not stable and the severity contrast is not identified. An $s=3$
switch rate above 80\% means the top level is a ceiling and severity is not graded; below 10\%
means the manipulation is too weak in this population. Opposite signs across models, or a
positive slope in fewer than two of three, fail the replication gate. Each of these is
pre-registered, so none of them can be reinterpreted after the fact as a discovery.



The limitations are known in advance and are part of the claim boundary. $\tau$-bench tasks are
public and may be memorized by target models, so baseline reliance could be inflated; the
within-prefix contrast remains valid, but the absolute switch rate may not transfer. The donor
prefixes come from gpt-4o and claude-3.5-sonnet, so the target models continue off-policy
histories; the design fixes prefixes across arms and stratifies by donor, but it does not make the
history on-policy. Severity is a constructed ordinal variable. If $s=1$ moves decisions as much as
$s=2$, the dose interpretation fails, and the paper reports that finding instead of relabelling
the levels.

The study covers three 8--9B model families, two domains, one renderer, deterministic tools, and a
four-step horizon. It does not claim end-to-end task success, frontier-model behaviour, live-API
validity, or a universal law. The tools are benchmark implementations against a pinned database,
not live services, so a wrong observation here is a controlled counterfactual rather than a
production incident. The annotation is done by a small author team, not an external vendor, which
limits the independence of the audit even with blinding. The 18~GB slice forces batch 4 rather
than the planned batch 8, and a mid-envelope G2 stop would reduce the sample and the precision of
the recovery estimate; both are disclosed if they occur. Finally, the manuscript is written
before the confirmatory data exist, so its Results section currently contains only the pilot and
the harvest counts, and every confirmatory statement is a placeholder with a named source.

\section{Conclusion}\label{sec:conclusion}

This paper asks whether the content of a real tool observation changes an agent's next decision,
whether the change is ordered by the severity of the content error, and how quickly the agent
recovers. The design holds 800 real $\tau$-bench prefixes fixed, varies only the observation across
five arms, and estimates a within-prefix conditional-logistic slope with a silent-propagation
decomposition and a four-step recovery curve. The study is pre-registered, gated, and runs at \$0
cash on a university H200 slice. The CPU-only harvest and freeze pipeline is implemented and the
sampling frame is measured; the human audit and the frozen manifest are pending, and no
confirmatory data have been collected. The contribution at this stage is the dataset, the
protocol, and the analysis plan. The empirical result, positive or null, will be reported at a
known precision and will not be reframed after the fact. \printcredits \section*{Declarations}
\begin{itemize}
  \setlength{\itemsep}{0pt}
  \setlength{\parsep}{0pt}
  \item \textbf{Funding:} This research received no external funding. The university provides the
  H200 MIG 1g.18gb slice in kind; the study has no cash budget and no paid API.   \item \textbf{Conflicts of Interest:} The authors declare that they have no known competing
  financial interests or personal relationships that could have appeared to influence the work
  reported in this paper.   \item \textbf{Data availability:} Raw $\tau$-bench trajectories are not redistributed. The
  upstream benchmark is MIT licensed and is pinned by commit and file hash in
  \texttt{data/manifests/source\_pin.json}. ToolResult-800 will release only scrubbed derived
  data: frozen prefix payloads, arm observations, perturbation records, renderer version, and
  hashes, under CC BY 4.0, with upstream notices preserved. Run ledgers stay local until scrubbed and
  packaged. The frozen manifest and the release checklist are the authoritative record.
  \pending{release DOI and manifest hash; the v1.0 artifact has not been released}.   \item \textbf{Code availability:} The harvest, replay, rendering, inference, coding, and
  analysis code is available at \url{https://github.com/Atik203/tool-observation} under the MIT
  license. Documentation and derived data are CC BY 4.0.   \item \textbf{Ethical approval:} The study collects no human-subject data. Annotation is
  performed by the author team on public benchmark rows, and direct identifiers are scrubbed
  before freeze. \pending{the institutional ethics determination for the annotation procedure}.   \item \textbf{Author contributions:} Conceptualization, data curation, formal analysis,
  investigation, methodology, project administration, software, validation, visualization, and
  writing (original draft) by M.A.R.; data curation, validation, and writing (review and editing)
  by S.A.S. and M.M.A.; supervision, validation, and writing (review and editing) by O.T. The
  CRediT allocation is carried over from the companion manuscript and must be confirmed by every
  author before submission.   \item \textbf{Use of AI-assisted tooling:} An AI-assisted mechanical pre-audit was used for the
  40-item G0 check under decision D-021, and its result is reported as provisional. The authors
  remain responsible for every label, every gate decision, and every claim in this paper.
  \pending{the journal-required AI-use statement}.
\end{itemize} \section*{Acknowledgements}
The authors thank the curators of $\tau$-bench for releasing the trajectories, task databases, and
harness that make a replayed fixed-prefix design possible, and the university high-performance
computing staff for the H200 MIG allocation. The authors also thank the members of the annotation
team who will complete the severity, outcome, and recovery audits. \hbadness=10000
\sloppy
\setlength{\bibsep}{0pt}
\renewcommand{\bibfont}{\fontsize{8.5pt}{9pt}\selectfont}
\bibliographystyle{elsarticle-num}
\bibliography{ref} \vspace{4pt}
\end{document}
